% ========================================================== INTRODUCTION ==
\section{Introduction}\label{sec:intro}

A traffic management centre does four things in sequence. It fills the gaps in its detector feed, it asks
where congestion will be in half an hour, it checks that its picture of the road is physically possible,
and it estimates where the traffic is coming from and going to. TrafficFlowBench \cite{tfb2026} turns that
loop into one leaderboard. Ten directional freeway panels, five-minute resolution and eleven months of
records are scored by
\begin{equation}
  \Stotal = 0.35\,\Sstate + 0.30\,\Squeue + 0.15\,\Sphys + 0.20\,\Sodme ,
  \label{eq:total}
\end{equation}
where each component lies in $[0,1]$. The tasks interlock by design. Physics is not submitted at all: it is read
off the reconstructed state. A reconstruction can match every detector to within noise and still create or
destroy vehicles between neighbouring links, a failure that squared error does not see. Queue forecasts begin
from the same congested region that the reconstruction must get right, and demand estimation needs counts that
are consistent with the flows.

The data are synthetic. A traffic-flow model whose parameters were measured from real freeway detector data
on these same corridors generated every split, so the records obey traffic physics and carry realistic measurement error, but no
external source contains the answers. Validation (March 2031) drives the public leaderboard and private
(April 2031) the final ranking. They are independently generated months with different incidents, five
in validation and seven in private, so a method tuned to the public month learns that month's incidents rather
than the traffic.

\subsection{What this report contributes}
\begin{itemize}
  \item \Tstate{} \textbf{A residual reconstruction model} that learns the error of same-day temporal interpolation
        from 92 features built only from the masked layer (\cref{sec:state}). It is reconciled with a predicted
        density, corrected on the free-flow branch by spatial covariance and low-rank factor models, averaged over
        three independently sampled fits, and finally recentred per link on short gaps where a two-stage bias test
        shows the same-day interpolant to be unbiased.
  \item \Tqueue{} \textbf{Queue forecasting under a strict twelve-slot information contract} (\cref{sec:queue}): forward-filled
        persistence that keeps dropped sensors, a joint-mask onset analogue that chooses the Bayes action for
        intersection-over-union over complete training outcomes \cite{nowozin2014}, and an ongoing-queue decoder that mixes
        a first-crossing classifier with a trajectory-analogue library \cite{lorenz1969}.
  \item \Todme{} \textbf{Exact split-specific demand estimation} by non-negative least squares
        \cite{lawson1974}, which reproduces the organizer's own reference solve (\cref{sec:odme}).
  \item \textbf{The leaderboard as an instrument} (\cref{sec:leaderboard}). The four tasks are scored on disjoint rows, so
        \cref{eq:total} is exactly additive. With the organizer baseline table as an anchor, single-component
        submissions identify component scores to rounding. We publish the full ledger of 31 submissions and the
        component values it implies.
  \item \Tphys{} \textbf{An honest account of what failed} (\cref{sec:physics,sec:negative}), including Godunov
        \cite{godunov1959,daganzo1994}, Newell \cite{newell1993} and diagram-inversion reconstructions that all lost to
        temporal interpolation on held-out months, and more than thirty queue, state and demand variants rejected by
        pre-registered gates.
  \item \textbf{Reproducibility engineering} (\cref{sec:repro}). TF3 and its predecessors passed an independent full-row
        parser and a clean regeneration of all 6,980,503 rows from the release and saved weights before upload; each
        file was uploaded exactly once and confirmed by a separate read-only listing.
\end{itemize}

\subsection{The two selected files}
The Kaggle dashboard selects two submissions (\cref{tab:selected}). TF3 is the cleanest end-to-end artefact:
it is regenerated from the raw release by one command, and that run matched all rows. ES4 is our best public
score. It keeps TF3's demand and flow, replaces the ongoing-queue labels with a classifier refit on all nine
training months, and removes a per-link speed bias on gaps of one to six slots. The two files contain the same
6,980,503 rows and differ in 5,825,667 speed values and all 87,000 ongoing-queue cells; onset labels, flows and
path flows are identical.

\begin{table}[h]
\centering\small
\caption{The two submissions selected for final scoring. Public scores are Kaggle results; no private score was
available at the time of writing.}\label{tab:selected}
\begin{tabular}{@{}l l l r l@{}}
\toprule
Submission & Kaggle ref. & Uploaded (UTC) & Public score & Archive SHA256 (first 16 hex) \\
\midrule
TF3 & 56825669 & 2026-10-04 13:51:11 & 0.86991 & \sha{2ae94d422399c724\ldots} \\
ES4 & 56998048 & 2026-10-09 07:13:22 & \textbf{0.87187} & \sha{ce03386b42e33bc6\ldots} \\
\bottomrule
\end{tabular}
\end{table}

\begin{notebox}[How to read the numbers in this report]
Only Kaggle returns leaderboard scores. Every other figure is a \emph{proxy}: an error measured on held-out
training months whose labels are released, on visible cells of the evaluation months that we masked ourselves,
or on the 80 published training queue windows. We say which proxy every number comes from. A proxy gain is not
a leaderboard forecast, and \cref{sec:transfer} shows how often our queue proxies failed to transfer.
\end{notebox}

\subsection{Organization}
\Cref{sec:benchmark} describes the benchmark and the properties of its design that shaped our method.
\Cref{sec:protocol} states the information boundaries and the experimental protocol. \Cref{sec:architecture}
gives the system architecture, and \cref{sec:state,sec:physics,sec:queue,sec:odme} treat the four tasks.
\Cref{sec:leaderboard} analyses the leaderboard evidence, \cref{sec:negative} lists negative results and
\cref{sec:discussion} discusses error, limitations and future work. \Cref{sec:repro} documents reproduction.
